%%%%%%%%%%%%%%%%%%%%%%%%%%%%%%%%%%%%%%%%%%%%%%%%%%%%%%%%%%%%%%%%%%%%%%%%

%%% LaTeX Template for AAMAS-2021 (based on sample-sigconf.tex)
%%% Prepared by Natasha Alechina and Ulle Endriss (version 2020-08-06)



%%%%%%%%%%%%%%%%%%%%%%%%%%%%%%%%%%%%%%%%%%%%%%%%%%%%%%%%%%%%%%%%%%%%%%%%

%%% Start your document with the \documentclass command.
%%% Use the first variant below for the final paper.
%%% Use the second variant below for submission.

\documentclass[sigconf]{aamas} 
%\documentclass[sigconf,anonymous]{aamas} 

%%% Load required packages here (note that many are included already).

\usepackage{balance} % for balancing columns on the final page
\usepackage[ruled,vlined]{algorithm2e}
\usepackage{subcaption}
\usepackage{array} % for aligning the items in the table as center
\usepackage{gensymb} % for the degree symbol
%%%%%%%%%%%%%%%%%%%%%%%%%%%%%%%%%%%%%%%%%%%%%%%%%%%%%%%%%%%%%%%%%%%%%%%%
\settopmatter{printacmref=false} % Removes citation information below abstract 
\renewcommand\footnotetextcopyrightpermission[1]{} % removes footnote with conference information in first column 
\pagestyle{plain} % removes running headers

%%% AAMAS-2021 copyright block (do not change!)

\setcopyright{ifaamas}
\acmConference[AAMAS '21]{Proc.\@ of the 20th International Conference on Autonomous Agents and Multiagent Systems (AAMAS 2021)}{May 3--7, 2021}{London, UK}{U.~Endriss, A.~Now\'{e}, F.~Dignum, A.~Lomuscio (eds.)}
\copyrightyear{2021}
\acmYear{2021}
\acmDOI{}
\acmPrice{}
\acmISBN{}

%%%%%%%%%%%%%%%%%%%%%%%%%%%%%%%%%%%%%%%%%%%%%%%%%%%%%%%%%%%%%%%%%%%%%%%%

%%% Use this command to specify your EasyChair submission number.
%%% In anonymous mode, it will be printed on the first page.

\acmSubmissionID{???}

%%% Use this command to specify the title of your paper.

\title{NovelGridworlds: A Benchmark Environment for Detecting\\ and Adapting to Novelties in Open Worlds}
%  \title{NovelGridworlds: Environments for Benchmarking Agents in Open World Novelties}

%%% Provide names, affiliations, and email addresses for all authors.

\author{Shivam Goel}
\authornote{The first two authors contributed equally.}
\affiliation{
%   \department{Department of Computer Science}
  \institution{Tufts University}
   \city{Medford, USA}}
\email{Shivam.Goel@tufts.edu}

\author{Gyan Tatiya}
\authornotemark[1]
\affiliation{
%   \department{Department of Computer Science}
  \institution{Tufts University}
  \city{Medford, USA}}
\email{Gyan.Tatiya@tufts.edu}

\author{Matthias Scheutz}
\affiliation{
  \institution{Tufts University}
  \city{Medford, USA}}
\email{Matthias.Scheutz@tufts.edu}

\author{Jivko Sinapov}
\affiliation{
%   \department{Department of Computer Science}
  \institution{Tufts University}
    \city{Medford, USA}}
\email{Jivko.Sinapov@tufts.edu}

%%% Use this environment to specify a short abstract for your paper.
\newcolumntype{P}[1]{>{\centering\arraybackslash}p{#1}} % for aligning the items in the table as center

\begin{abstract}
% This document outlines the formatting instructions for submissions to
% AAMAS-2021. You can use its source file as a template when writing 
% your own paper. It is based on the file `\texttt{sample-sigconf.tex}'
% distributed with the ACM article template for \LaTeX\@.
As researchers are developing methods for detecting and accommodating novelties that will make AI agents more robust to unknown sudden changes in the ``open worlds'', there is an increasing need for benchmark environments that allow for the systematic evaluations of the proposed AI techniques.
% Moreover, the need of benchmarking environment to develop these agents is necessary since the present benchmarking environments ...
We present ``NovelGridworlds'', an {\em OpenAI Gym} environment framework for developing and evaluating AI agents that can detect and adapt to unknown sudden novelties in their environments.  Based on a rich taxonomy of novelties illustrated in $4$ different tasks with $12$ novelties in each, we propose evaluation metrics for evaluating both planning systems and learning systems that can handle novelty and illustrate the metrics with results from simulations of both types of AI agents.

% We also demonstrate different artificial intelligence techniques to serve as a benchmark for researchers to extend their work, namely planning methods and reinforcement learning methods
\end{abstract}

%%% The code below was generated by the tool at http://dl.acm.org/ccs.cfm.
%%% Please replace this example with code appropriate for your own paper.


%%% Use this command to specify a few keywords describing your work.
%%% Keywords should be separated by commas.

\keywords{Learning and Adaptation, Novelty, Non-stationary MDP, Reinforcement Learning, Planning, Benchmarking Environments}

%%%%%%%%%%%%%%%%%%%%%%%%%%%%%%%%%%%%%%%%%%%%%%%%%%%%%%%%%%%%%%%%%%%%%%%%

%%% Include any author-defined commands here.
         
\newcommand{\BibTeX}{\rm B\kern-.05em{\sc i\kern-.025em b}\kern-.08em\TeX}

%%%%%%%%%%%%%%%%%%%%%%%%%%%%%%%%%%%%%%%%%%%%%%%%%%%%%%%%%%%%%%%%%%%%%%%%

\begin{document}

%%% The following commands remove the headers in your paper. For final 
%%% papers, these will be inserted during the pagination process.

\pagestyle{fancy}
\fancyhead{}

%%% The next command prints the information defined in the preamble.

\maketitle 

% %%%%%%%%%%%%%%%%%%%%%%%%%%%%%%%%%%%%%%%%%%%%%%%%%%%%%%%%%%%%%%%%%%%%%%%%

\section{Introduction}

% In developing artificial intelligent agents to sustain in real world, it is pertinent for the agents to be habitual to changes in their operating environments. Like in real world, the environment can suddenly change and humans tend to adapt to them quickly, the AI agents also needs to be able to do the same. In an attempt to make the algorithms scalable and generalizable to real world domains, researchers are developing Reinforcement learning and planning algorithms to deal with open world novel scenarios. 
%Human behavior is highly flexible, able to persevere through unexpected events, and generalize to %unfamiliar tasks. Likewise, 
Artificial agents operating in real-world applications need the ability to be robust to unknown sudden changes in their operating environment.  Yet, most research in AI has been based on ``closed-world assumptions'', i.e., that all task-relevant information, at least the involved types of objects and the types of attributes they have, is known ahead of task performance.  In contrast, ``open worlds'' may contain novelties that are unknown during algorithm design and can, therefore, not be accounted for ahead of time. There are increasing efforts for developing AI agents to adapt to novelties in open worlds, using techniques such as reinforcement learning (RL)~\citep{khetarpal2020towards, padakandla2020reinforcement}, planning-based approaches~\citep{muhammad2021novelty, klenkmodel}, and hybrid approaches~\citep{sarathy2020spotter, peng2021detecting}.

When developing algorithms for open worlds, researchers need to use domains to train and evaluate their agents, and the domains are likely tailored to the agent's algorithms and thus do not allow for a fair comparison across different approaches.  What is needed is an evaluation environments capable of benchmarking different types of AI agents without introducing a bias towards one or another method. Moreover, evaluation environments with cross-platform support for different programming languages and support for various sub-domains of AI are not readily available.  To address this problem, we present a ``gym environment'' called {\em NovelGridworlds}\footnote{The open-source NovelGridworlds repository is available at~\url{https://github.com/gtatiya/gym-novel-gridworlds}.}

% - we want to talk about novelty characterization and why its important
% - we also want to talk about why our domain is a better domain than others.
% - we also want to 
The rest of the paper is organized as follows. We start by providing background on the existing environments used for benchmarking such agents in Section~\ref{sec:related_works}, followed by providing basics of the planning agents and the reinforcement learning agents in Section~\ref{sec:prelims}. We then describe the base tasks in the NovelGridworlds domain in Section~\ref{sec:environment}, followed by a detailed taxonomy of novelty with examples of the same implemented in NovelGridworlds. In Section~\ref{sec:experiment}, we present the agents used with our domain, followed by the proposed evaluation metrics, and the results in Section~\ref{sec:results}.  We finish with a discussion of our accomplishments and possible future improvements in Section~\ref{sec:discussion} and Section~\ref{sec:conclusion}, respectively.

Our contributions are thus (1) a benchmarking environment for evaluating and prototyping multiple AI agents that can detect and handle novelties; (2) a rich taxonomy of novelties and illustration of $4$ tasks with $12$ novelties in each of the tasks; and (3) evaluation metrics to evaluate both planning agents as well as reinforcement learning agents with baseline results using both the agents on some tasks and novelties. 



% - Motivate the idea of the environment. \\
% - Easy to prototype the novel scenarios.\\
% - Used in both planning and learning domains. \\
% - Light weight, 2-D representation of the environments. \\
% - Mention Tom Silver limitations. \\
% - This is not deterministic and has stochasticity to model novelty. \\
% - New tasks can be easily created by a simple pull request. \\

% Reinforcement Learning (RL) agents face enormous difficulties in learning complex tasks. However, this difficulty is significantly increased if the tasks tend to change their dynamics. Many methods employ the use of lifelong learning to address these issues~\citep{tessler2016deep}. These methods formulate the problem as a continual learning scenario, and adapt to learn the changes that occur over time~\citep{xie2020deep}.

% The methods dealing with environment changes are still less, and even the ones we currently have, suffer from two main limitations. First, the methods do not deal with complex environments, specifically environments that have sub goals and a structure. As a result, if the same frame work is applied to complex tasks like Minecraft, the existing framework is weak to handle the complexity with adaption. Second, the changes in the environment are assumed to be gradual. If there is a sudden random change in the environment that can hamper the learning ability these methods become brittle and some take long training times to learn the changes. 

% To address these shortcomings, this work presents a hierarchical RL framework which has the capability to learn complex tasks as well as has the capability to accommodate changes in a few training steps.
% We model the architecture of our agent using Hierarchical Reinforcement Learning (HRL). This architecture has the capability of learning on the fly when a new entity appears in the environment or a state variable changes. We call the changes occurring in the environment as novelties. Novelty detection, however, is not the focus of this paper and thus we assume that the novelty is already detected by a system. Hence, the proposed system starts off after the onset of novelty.

% We evaluate our method by producing novelties on a grid-world similar to minecraft. The novelties are generated at various levels of difficulties and we evaluate our approach on three novel environments. Our experiments show that our architecture can adapt to these changes and has the capability of performing learning on the fly in a significantly less training steps.

% In this paper our contributions are as follows:
% \begin{itemize}
%     \item A Hierarchical Reinforcement learning (HRL) framework which uses symbolic state representation as input and solves extremely complex tasks such as minecraft without any human demonstrations.
%     \item A HRL framework to perform learning on the fly.
%     \item Empirical evaluations on a custom grid world domain in which we add sudden changes to the environment in the form of $3$ different novel scenarios. We show that the proposed algorithm is capable of performing learning on the fly while training on significantly less time steps.
% \end{itemize}
% . TODO: {Add citations related to the dynamically changing methods}. These frameworks are efficient at solving tasks which are relatively easier and the environment dynamics changes gradually TODO: {cite here}. Our work focuses on tasks that are more challenging than Pacman or Mario, and have the capability of catastrophically changing environmental dynamics. These changes in the environment dynamics can impact the agents' learning performance drastically. In this work we consider a highly complex environment like Minecraft and pose the challenge of a environmental dynamics that changes suddenly and can lead to performance issues of the learning algorithm. 
% xThis document explains the main features of the `\texttt{aamas}' 
% document class, which is essentially identical to the `\texttt{acmart}'
% document class provided by the ACM. The only difference is a minor 
% modification to allow for the correct copyright attribution to IFAAMAS.
% For detailed documentation of the original document class, please refer
% to the relevant website maintained by the~ACM:
% %
% \begin{center}
% \url{https://www.acm.org/publications/proceedings-template}
% \end{center}
% %
% The first command in your source file should be either one of these:
% \begin{verbatim}
%     \documentclass[sigconf]{aamas}
%     \documentclass[sigconf,anonymous]{aamas}
% \end{verbatim}
% %
% The first variant should be used for final papers. The second should be
% used when you submit your paper for blind review; it will replace the 
% names of the authors with the submission number.

% Make sure your paper includes the correct copyright information and 
% the correct specification of the \emph{ACM Reference Format}. Both of 
% these will be generated automatically if you include the correct 
% \emph{copyright block} as shown in the source file of this document.

% Modifying the template---e.g., by changing margins, typeface sizes, 
% line spacing, paragraph or list definitions---or making excessive use 
% of the `\verb|\vspace|' command to manually adjust the vertical spacing 
% between elements of your work is not allowed. You risk getting your 
% submission rejected (or your final paper excluded from the proceedings) 
% in case such modifications are discovered. The `\texttt{aamas}' document 
% class requires the use of the \textit{Libertine} typeface family, which 
% should be included with your \LaTeX\ installation. Please do not use 
% other typefaces instead.

% Please consult the \emph{Call for Papers} for information on matters 
% such as the page limit or anonymity requirements. It is available from
% the conference website:
% %
% \begin{center}
% \url{https://aamas2021.soton.ac.uk/}
% \end{center}
% %
% To balance the columns on the final page of your paper, use the 
% `\texttt{balance}' package and issue the `\verb|\balance|' command
%  somewhere in the text of what would be the first column of the last 
%  page without balanced columns. This will be required for final papers.

%%%%%%%%%%%%%%%%%%%%%%%%%%%%%%%%%%%%%%%%%%%%%%%%%%%%%%%%%%%%%%%%%%%%%%%%
\section{Background}
\subsection{Related Work}
\label{sec:related_works}
In the recent past, there is a growing interest in the AI community to develop agents that can adapt to changes in open-world environments.
Khetarpal et al. ~\citep{khetarpal2020towards}, presents a review of methods used to adapt to non-stationarity in environments using reinforcement learning methods.
Muhammad et al. ~\citep{muhammad2021novelty} propose a cognitive architecture using open world symbolic planning to solve novelties in an open-world environment. 
Jamshidi et al.~\citep{jamshidi2019machine} present an architecture that utilizes machine learning techniques to find  Pareto-optimal configurations and limit the vast search space to such configurations to make planning feasible.
Klenk et al. ~\citep{klenkmodel} propose HYDRA, which uses a domain-independent planner to play the game, and learns a model to update it when novelties get introduced in the game. Sarathy et al. ~\citep{sarathy2020spotter} propose SPOTTER, which uses reinforcement learning techniques to find the operators to solve the tasks when symbolic planning fails to solve them.
There has been an enormous reliance on the choice of domains to experiment upon in all of these works. Hence, it becomes necessary to develop domains best suited for the problems at hand. 

The development of environments and benchmarks has helped progress the research in various fields of AI~\citep{wang2018glue, fu2020d4rl, nichol2018gotta, pardo2020tonic, silver2020pddlgym}. Alex et al. ~\citep{wang2018glue} present GLUE, a benchmarking environment platform to bridge learning with natural language understanding. Jiang et al. ~\citep{jiang2020wordcraft} present WordCraft, an environment to develop common-sense agents. This environment also incorporates natural language corpora to improve common sense among agents.
Nichol et al.~\citep{nichol2018gotta}, present a benchmarking environment to evaluate the effectiveness of transfer learning and few-shot learning techniques in the context of Reinforcement Learning. To bridge the gap for the development of algorithms with a combination of planning and learning, PDDLGym~\citep{silver2020pddlgym} offers an environment to combine PDDL (Planning Domain Definition Language) with OpenAI Gym~\citep{1606.01540} environments.

However, environments with cross-platform support for various programming languages and support for various sub-domains of AI are not that easy to find. Moreover, NovelGridworlds provide tools to inject several novelties in all the base tasks. To best of our knowledge there is no such environment that provides novelty injection tool. 
% Thus, we present NovelGridworlds, an OpenAI Gym~\citep{1606.01540} environment that offers cross-platform support, and we also show the same environment can be used for both planning systems as well as learning systems.

\subsection{Notations \& Preliminaries}
\label{sec:prelims}
In sequential decision making problems, the agent has to perform decision making at each time step to reach a set of goal states $g$ from a start state $s_o$. However, in solving the problems in open-worlds, there can be unexpected changes in the environments, which can deter the agent from reaching the goal state. As elucidated in~\citep{muhammad2021novelty}, we borrow the definition of novelty as something that the agent has not experienced before, nor it can derive it using some known experience. 
We define these unexpected changes as novelties. Ideally, the agents should have the capabilities to still reach the goal state even if novelties arise in the environments. In this paper, we extend this definition by providing a taxonomy of the classes of novelties and present examples of each of them implemented in our NovelGridworlds domain.

% The methods that AI agents use to perform decision making can be broadly classified into two categories, planning and learning. We now define how a planning and RL problem is formulated and their similarities and differences. A typical planning problem involves decision making to reach from a start state to a goal state. We define a planning problem using a tuple $\Sigma=(\mathcal{S}, \mathcal{A}, \tau)$, where $\mathcal{S}$ is a set of finite states of the system, $\mathcal{A}$ is a finite set of actions, and $\tau$ representing the transition function between the states and the actions. In a deterministic world, the planning agent is given information of the states and actions. The agent searches for a plan to reach the goal state. 
\subsubsection{Planning}
To define a planning problem, we define $\mathcal{L}$, as a first-order language that contains atoms $p\left(t_{1}, \ldots, t_{n}\right)$ and their negations $\neg p\left(t_{1},\ldots, t_{n}\right)$ , where $t_{i}$ can be either variables or constants. We define a planning domain in the language $\mathcal{L}$ as $\Sigma=(\mathcal{S}, \mathcal{A}, \tau)$, in which $\mathcal{S}$ represents the set of states, $\mathcal{A}$ represents the set of finite actions, and $\tau$ is the transition function between the states and actions~\citep{sarathy2018macgyver}. A typical planning problem is defined as $\mathcal{P}=\left(\Sigma, S_{0}, G\right)$, in which $S_{0},$ is the set of starting states and $G$ is the set of goal states.
The agent begins in one of start states $S_{0}$, and finds a plan $\pi$ to reach one of the goal state in $G$. Thus the plan $\pi=\left[a_{1}, a_{2}, \ldots a_{|\pi|}\right]$ is the solution to the planning problem $\mathcal{P}$ where $a_{i} \in \mathcal{A}$ is the action.
Each action has a set of preconditions and effects. The set of preconditions needs to be true for the action to be executed successfully in the environment. Table~\ref{tab:preconds} shows examples of how actions are defined in planning domain. Note that these actions are only for demonstration and can be defined using a different set of pre-conditions and effects.

% The set of states are composed of objects in the environment and objects have fluents and affordances.

% in planning, one knows the effects of one's actions, i.e., the transition function is known to the agent,
% - in planning, a known precondition-condition function pre: $S \times A \rightarrow\{$ true, false $\}$ is given, which specifies in which states it is legal to apply which actions.
% - in planning, one is given a goal function goal : $S \rightarrow$ true, false $\},$ which characterizes the target states.
% - in planning, the aim is to start from a state $s_{1}$ and to find a sequence of actions $a_{1}, \ldots, a_{n}$ $\left(a_{i} \in A\right)$ such that
% o goal $\left.\left(\left(\ldots \quad\left(s, a_{1}\right)\right) \ldots, a_{n-1}\right), a_{n}\right)=$ true, and
% - pre( (... $\left.\left.\left(s, a_{1}\right)\right) \ldots, a_{i-1}$ ), $\left.a_{i}\right)=$ true.
\subsubsection{Reinforcement learning}
A RL problem is generally formalized using a Markov Decision Process (MDP) that is represented by a tuple $\langle S, A, \tau, R, \gamma\rangle .$ At every time step $t,$ the $\mathrm{RL}$ agent receives a state representation $S_{t} \in \mathcal{S}$ and explores an unknown environment by taking an action $A_{t} \in \mathcal{A}(s) .$ A reward $R_{t+1} \in \mathcal{R} \subset \mathbb{R}$ is provided based on the action taken by the agent to reach the next state $S_{t+1}$. The agent learns to maximize the expected return value $G_{t}=\sum_{k=0}^{\infty} \gamma^{k} R_{t+k+1}$ for every state at time $t .$ The discount factor $\gamma \in[0,1)$ determines the importance of immediate and future rewards ~\citep{sutton2018reinforcement}.
% - Talk about the planning and rl agent commonality in the sense of A MDP tuple. \\
% - Formalize the specifics of the MDP to develop the idea of novelty classification later on. \\
% %
% \begin{figure}[h]
%   \centering
%   \includegraphics[width=0.6\linewidth]{gridworld.png}
%   \caption{The figure shows a diagrammatic representation of the novelgridworld domain. A wall surrounds the entire arena of the gridworld, and the agent is facing a crafting-table. The space in green is occupied by air, and the agent can move through that space to navigate in the arena.}
%   \label{fig:arena}
% \end{figure}
% %

%
\begin{table}[!tb]
  \centering
  \includegraphics[width=0.9\linewidth]{preconds_effects_4.png}
  \caption{Examples of the \emph{break} and \emph{move-forward} action as represented in the planning domain.}
  \label{tab:preconds}
\end{table}
%

\section{Environment}
\label{sec:environment}
NovelGridworlds is an OpenAI Gym environment which is implemented in python and currently has four different tasks.
% \begin{itemize}
%     \item NovelGridWorld-bow-$v0$
%     \item NovelGridWorld-bow-$v1$
%     \item NovelGridWorld-pogostick-$v0$
%     \item NovelGridWorld-pogostick-$v1$
% \end{itemize}
Each environment is a discrete grid world domain. The size of the default grid world is set to $10$, which the user can vary. Each block contains an entity or object. The contents of the blocks can vary according to the tasks. However, in each environment, there are a few common block items. The entire arena is enclosed by a \emph{wall}, which is unbreakable. Each environment has a block occupied by a \emph{crafting- table}, a block occupied by the \emph{agent}, and free space occupied by \emph{air}. The agent can move in the blocks containing air. At each time step, the agent maintains an inventory of the collected items. A diagrammatic representation of the grid world is shown in figure~\ref{fig:arena}. The diagrammatic representation is not specific to a particular task. It only shows the \emph{agent}, \emph{wall}, \emph{crafting-table} and \emph{air}. It also demonstrates how a novel entity appears in the environment.  
The agent is given a discrete set of actions, out of which $4$ actions are common to all the environments, namely \emph{move-forward}, \emph{turn-left}, \emph{turn-right} and \emph{break}. Some of the grid world blocks are occupied by breakable objects. The agent when one block next to these objects uses \emph{break} action, the item falls into agent's inventory. For example, if the agent uses \emph{break} action in front of the \emph{tree-log}, the \emph{tree-log} falls into the agent's inventory.
% For the sake of simplicity, \emph{crafting-table} is made unbreakable object while training RL agents.
The environment also returns an action cost for each action taken by the agent. This is designed by keeping in mind a planning agent to optimize the plan based on the action cost.
In each task, the agent needs to craft objects using multiple recipes, which requires performing certain steps in some sequence.
We describe each task in detail in the sub-sections below.

% % %
% \begin{table}
% \begin{center}
% \begin{tabular}{ |P{3cm}||P{2cm}| }
%  \hline
%  \multicolumn{2}{|c|}{\textbf{NovelGridworld-Bow-v0}} \\
%  \hline
%  \textbf{Predicates}& physical, direction,world 
%  \hline
 
% Preconditions   & (and 
%         (at agent loc01) 
%         (orientation agent north) 
%         (adjacent loc01 loc02 north) 
%         (at tree-log loc02) 
%         (not (unbreakable tree-log))
%     ) \\
% Effects   & $80.00 \pm 40.00$ \\
%   Fire wall   & $0.00 \pm 0.00$ \\
%     Break increase & $29.2 \pm 19.6$ \\
%  \hline
% \end{tabular}
% \caption{\label{tab:results_rl}Table shows the results of the performance of the reinforcement learning agent on the base task of \emph{NovelGridworld-Bow-v0} and three novelty scenarios. The values denote the mean and standard deviations of the success rate of solving the final task. The values were calculated at 1.8M timestep for the default case and 4 M time step for all the novelty scenarios.}
% \end{center}
% \end{table}
% %
%
\begin{figure}[!tb]
  \centering
  \includegraphics[width=1\linewidth]{screen-shot_ng.png}
  \caption{The figure shows a screenshot of one of the states of the NovelGridworld-Bow-v0 of size $10\times10$. The agent is represented as a triangle, and it is facing the crafting table (shown in green). Two of the three sticks (yellow) and strings (orange) are in gridworld and one of each is in the agent's inventory. The spaces in blue are the air blocks, and the agent can freely traverse in those spaces. The entire arena is surrounded by walls represented in purple. The image's left side shows the steps, agent's current direction, last action taken, reward, step cost, and selected item. On the right side of the image, the agent's inventory is displayed.}
  \label{fig:screenshot_ng}
\end{figure}
%
\subsection{Task Descriptions}
\subsubsection{\textbf{NovelGridworld-Bow-v0}}
\label{sec:bow-v-0}
In this task, the agent has to achieve the goal of successfully crafting a \emph{bow}. Apart from the \emph{crafting-table}, \emph{wall}, \emph{air}, and \emph{agent}, three blocks in the grid world are occupied by \emph{string} and three blocks are occupied by \emph{stick}. In addition to the default actions common to all the domains (described earlier), this task has an additional action called \emph{craft-bow}. The \emph{recipe} for creating a bow requires three \emph{stick} and three \emph{string}. After collecting all the required ingredients from the environment, the agent needs to navigate to the \emph{crafting-table} to craft a \emph{bow}. Upon successful execution, the \emph{bow} appears in the agent's inventory, and the task is complete. A screenshot of this environment can be seen in figure~\ref{fig:screenshot_ng}.


%
\begin{figure}[tb]
  \centering
   \begin{subfigure}[c]{0.70\linewidth}
      \includegraphics[width=\linewidth]{gridworld.png}
      \caption{Default (No novelty)}
  \end{subfigure}
  \begin{subfigure}[c]{0.30\linewidth}
    \includegraphics[width=\linewidth]{gridworld_axe.png}
     \caption{Axe}
  \end{subfigure}
  \begin{subfigure}[c]{0.30\linewidth}
    \includegraphics[width=\linewidth]{gridworld_fence.png}
    \caption{Fence}
  \end{subfigure}
  \begin{subfigure}[c]{0.30\linewidth}
    \includegraphics[width=\linewidth]{gridworld_firewall4.png}
    \caption{Fire-wall}
  \end{subfigure}
  \caption [The figures show a diagrammatic representation of the NovelGridworlds domain and representation of some of the novelties. From subfigure (a), it can be seen that a wall surrounds the entire arena of the grid world, and the agent is facing a crafting-table. The green space represents the air, and the agent can move through that space to navigate inside the arena. In subfigure (b), a new object (axe) appears in the environment. In subfigure (c), a new object (fence) surrounds the crafting table to obstruct it from the agent. In subfigure (d), the walls are replaced by fire-walls.]
  {The Figures show a diagrammatic representation~\footnote
  {To reduce clutter, we only show the default items in the environment. Hence, the diagrammatic representation only shows the \emph{agent}, \emph{crafting-table}, \emph{wall}, and \emph{air}} of the NovelGridworlds domain and representation of some of the novelties. From subfigure (a), it can be seen that a wall surrounds the entire arena of the grid world, and the agent is facing a crafting-table. The green space represents the air, and the agent can move through that space to navigate inside the arena. In subfigure (b), a new object (axe) appears in the environment. In subfigure (c), a new object (fence) surrounds the crafting table to obstruct it from the agent. In subfigure (d), the walls are replaced by fire-walls.}
  \label{fig:arena}
\end{figure}
%
\footnotetext{To reduce clutter, we only show the default items in the environment. Hence, the diagrammatic representation only shows the \emph{agent}, \emph{crafting-table}, \emph{wall}, and \emph{air}.}
\subsubsection{\textbf{NovelGridworld-Bow-v1}}
\label{sec:bow-v-1}
As opposed to the previous task (\ref{sec:bow-v-0}), in this task, the agent has to collect raw ingredients from the environment and craft the intermediaries to craft the \emph{bow}. This task is an advanced version of the previous task. In addition to the default grid world items, the blocks in this grid world are occupied by \emph{tree-log} and \emph{wool}. The default action space is augmented by \emph{extract-string}, \emph{craft-plank}, \emph{craft-stick} and \emph{craft-bow}. The agent can use the action ~\emph{craft-plank} to convert $1$ \emph{tree-log} into $4$ \emph{plank}, and use \emph{craft-stick} to make $4$ \emph{stick} from $2$ \emph{planks}. The agent can use \emph{extract-string} in front of \emph{wool} to extract $4$ \emph{string} from the \emph{wool}. The task's final goal is to craft a \emph{bow}, and the agent needs to be in front of the \emph{crafting-table} with $3$ \emph{stick} and $3$ \emph{string} to complete the task.
\subsubsection{\textbf{NovelGridworld-Pogostick-v0}}
\label{sec:pogo-v-0}
In this task, the agent has to craft a \emph{pogostick}. The grid world in this task has \emph{tree-log}, \emph{tree-tap}, \emph{stick}, \emph{plank}, in addition to the default items. The default action space is augmented by \emph{extract-rubber} and \emph{craft-pogostick}. One of the \emph{tree-log} has the \emph{tree-tap} placed in front of it. The agent needs to extract \emph{rubber} from that unique \emph{tree-log}. The agent needs $4$ \emph{stick}, $2$ \emph{plank}, and $1$ \emph{rubber} to craft a \emph{pogostick}. The agent needs to be facing the \emph{crafting-table} to craft the \emph{pogostick}. The task is complete when a \emph{pogostick} gets added to the agent's inventory.
\subsubsection{\textbf{NovelGridworld-Pogostick-v1}}
\label{sec:pogo-v-1}
It is an advanced version of the task in~\ref{sec:pogo-v-0}. In addition to the default items, the grid world in this task has some blocks occupied with \emph{tree-log}. The agent's default action space is expanded by \emph{place-tree-tap}, \emph{extract-rubber}, 
\emph{craft-plank}, \emph{craft-stick}, \emph{craft-tree-tap}, and, \emph{craft-pogostick}. In order to complete the task of successfully crafting the \emph{pogostick}, the agent needs to craft intermediaries. The agent needs $1$ \emph{tree-log} to craft $4$ \emph{plank}, $2$ \emph{plank} to craft $4$ \emph{stick}, and $5$ \emph{plank} and $1$ \emph{stick} to craft a \emph{tree-tap}. The \emph{tree-tap} can be placed in any of the \emph{tree-log} to extract \emph{rubber}. The recipe to craft the \emph{pogostick} requires $4$ \emph{stick}, $2$ \emph{plank}, and $1$ \emph{rubber}. To craft the \emph{tree-tap} and \emph{pogostick}, the agent needs to be facing the \emph{crafting-table} with the necessary ingredients. Once the agent gets the \emph{pogostick} in its inventory, the task is complete. Table~\ref{tab:actions} summarizes the actions and blocks in each task.


\subsection{Novelty Taxonomy}
As described in Section ~\ref{sec:prelims}, a planning agent and a RL agent both have access to a finite set of states $\mathcal{S}$, set of finite actions $\mathcal{A}$, and only exclusive to a planning agent, access to the transition function $\tau$. Any change in these sets would be considered novel to the AI agent. Therefore, to cover extensive novel scenarios, we classify the novelties by changing any or all of these sets. As described below, we divide the novelties into three broad classes.



\subsubsection{Object Novelty}
In this class of novelty, a novel object appears in the environment. This object is completely unknown to the agent's knowledge of the world. In other words, the observation space $x \in \mathcal{S}$ does not come from the distribution of the agent's current perception of the world, and hence in this class of novelty, state $\mathcal{S}$ changes.
To illustrate an example, let us consider the \emph{NovelGridworld-Pogostick-v1} task as described in Section ~\ref{sec:pogo-v-1}. In the post-novelty scenario, a novel object called \emph{axe} appears in the environment. The agent's current perception of the world has no notion of \emph{axe} and hence it is considered as a novelty. The agent has to learn to interact with the \emph{axe} to figure out its usage in goal completion.
\subsubsection{Attribute Novelty}
In this class of novelty, the attributes of the existing items in the environment changes. We define this class of novelty using the definition of transition function $\mathcal{\tau}: \mathcal{S} \times \mathcal{A} \rightarrow \mathcal{S}$, as described in Section ~\ref{sec:prelims}. In the post-novelty scenario, the environment changes such that $\mathcal{\tau^{\prime}}: \mathcal{S} \times \mathcal{A} \rightarrow \mathcal{S^{\prime}}$ in which $\mathcal{S^{\prime}}$ is the new state transition from the state action pair. Hence, similar state action pairs do not result in the same states as in the pre-novelty scenario. To illustrate this novelty, let us take an example of the \emph{NovelGridworld-Pogostick-v1} task as described in Section~\ref{sec:pogo-v-1}. In the post-novelty scenario, if the \emph{tree-log} provides $2$~\emph{tree-log} into the agent's inventory as opposed to $1$, then we call it a attribute change of the \emph{tree-log}. Therefore, the agent encounters a new state $\mathcal{S^{\prime}}$ from known state $\mathcal{S}$ taking known action $\mathcal{A}$.
\subsubsection{Action Novelty}
In this class of novelty, the action space $\mathcal{A}$ of the environment changes. In other words, the behavior of action as known to the agent is not the same in the post-novelty scenario. To illustrate this class of novelty let us consider the agent's action space to be a set of three actions, namely, \emph{move-forward} which takes the agent one step forward in the direction the agent is facing; \emph{turn-left}, which turns the agent $90\degree$ counter-clockwise; and \emph{turn-right}, which turns the agent $90\degree$ clockwise. In the post-novelty scenario, if the agent moves one block backward on using the action \emph{move-forward} or if the agent turns $90\degree$ clockwise on using the action \emph{turn-left}, then the agent's perception of the actions will change and hence, we call this as action novelty. We provide specific cases of each class of novelty later in the Section~\ref{sec:novel_domains}
% \begin{itemize}
%     \item Object Novelty: 
%     \item Attribute Novelty
%     \item Action Novelty
% \end{itemize}
% \emph{Object Novelty} pertains to the novelties in which a novel object appears in the environment. This novel object is unknown to the trained agent. This novelty class covers the aspect of the state space $\mathcal{S}$ changes of the environment.
% \emph{Attribute Novelty} is the class of novelty that deals with the changes in the $\mathcal{\tau}$ of the agent's environment. In this novelty, the attributes of the existing items in the environment change, resulting in a change in the environment's transition dynamics.
% \emph{Action Novelty}, is the class of novelty in which the changes in the action space $\mathcal{A}$ take place. We use two examples to explain this. In one of the novel instances, the action space of the environment is jumbled up. Another instance of this class of novelty can be an addition of an entirely new action.

We further divide each class of novelty into three sub types, \emph{Beneficial}, \emph{Detrimental}, and \emph{Irrelevant}.

\subsubsection{Beneficial}
This type of novelty is helpful for the agent to complete the task at hand. Hence, the agent needs to learn to utilize this novelty in task completion.

\subsubsection{Detrimental}
This type of novelty is further catastrophic in task completion. In other words, the agent needs to learn to mitigate this novelty to complete the task efficiently.

\subsubsection{Irrelevant}
This type of novelty is irrelevant to the task. In other words, the novelty can be completely ignored by the agent.


%
\begin{table}
\begin{center}
\begin{tabular}{ |P{2.5cm}||P{2.5cm}| |P{1.6cm}|}
 \hline
 \multicolumn{3}{|c|}{\textbf{Actions and Blocks in Tasks}} \\
 \hline
 \textbf{Environment}& \textbf{Actions} & \textbf{Blocks} \\
 \hline
 Default   & \small{{Move-forward, turn-left, turn-right, break, select-<block>.}} &\small{{air, wall, crafting-table}} \\
 \hline
  NovelGridworld-Bow-v0   & \small{craft-bow.} & \small{{string, stick}} \\
  \hline
   NovelGridworld-Bow-v1   & \small{{craft-bow, extract-string, craft-plank, craft-stick.}} & \small{{tree-log, wool, string, stick}} \\
   \hline
    NovelGridworld-Pogostick-v0 & \small{{extract-rubber, craft-pogostick.}} & \small{{tree-log, tree-tap, stick, plank}}\\
    \hline
       NovelGridworld-Pogostick-v1 & \small{{extract-rubber, craft-plank, craft-stick, craft-tree-tap, craft-pogostick.}} &\small{{tree-log, tree-tap, stick, plank}} \\
 \hline
\end{tabular}
\caption{\label{tab:actions}Table shows the actions and blocks specific to each of the implemented environments. Note that "Default" denotes the actions and blocks common to all the environments.}
\end{center}
\end{table}
%

% \begin{itemize}
%     \item Beneficial
%     \item Detrimental
%     \item Irrelevant
% \end{itemize}
% In \emph{Beneficial}, 
% However in \emph{Detrimental},  In \emph{Irrelevant}, 
% The new entity can serve three purposes, either it can be useful to the agent in the task completion. Hence the agent needs to learn to exploit it to perform better. It can also be detrimental to the current task.

% You will be assigned a submission number when you register the abstract 
% of your paper on \textit{EasyChair}. Include this number in your 
% document using the `\verb|\acmSubmissionID|' command.\\
% - Dynamically changing Environments \\

% - Hierarchical Reinforcement learning \\

% Then use the familiar commands to specify the title and authors of your
% paper in the preamble of the document. The title should be appropriately 
% capitalised (meaning that every `important' word in the title should 
% start with a capital letter). For the final version of your paper, make 
% sure to specify the affiliation and email address of each author using 
% the appropriate commands. Specify an affiliation and email address 
% separately for each author, even if two authors share the same 
% affiliation. You can specify more than one affiliation for an author by 
% using a separate `\verb|\affiliation|' command for each affiliation.

% Provide a short abstract using the `\texttt{abstract}' environment.
 
% Finally, specify a small number of keywords characterizing your work, 
% using the `\verb|\keywords|' command. 

%%%%%%%%%%%%%%%%%%%%%%%%%%%%%%%%%%%%%%%%%%%%%%%%%%%%%%%%%%%%%%%%%%%%%%%%

% \section{Learning Agents}
% - Describe the two learning agents and their capability with respect to the environment in detail.

% \subsection{Planning Agent}
% -Describe the planning from PDDL representation\\
% -Possibly add the automatic generation of PDDLs from the sense-all commands\\

% \subsection{Reinforcement Learning Agent}
% - Describe the MDP details of the RL agent.\\
% - State space: Observation\\
% - reward function and action spaces\\
% - Feature representation of the state space.\\
% - HRL framework description. \\
% % For help with typesetting the body of your paper in \LaTeX\@, please 
% % make use of the familiar resources~\cite{Lam94}. In this section we 
% % % merely highlight a few specific features. 
% % In this section, we begin by describing the hierarchical reinforcement learning architecture used to train the base task. We then describe the proposed method to adapt to the changes in the environment. We also describe the evaluation metrics used to evaluate the proposed method.
%
\begin{figure}[!tb]
  \centering
  \includegraphics[width=1.1\linewidth]{venn_diag_4.png}
  \caption{Classification of the novelties implemented in the NovelGridworlds. The novelties in this diagram are implemented across all the domains presented in the paper.}
  \label{fig:venn_diag}
\end{figure}
%
%
\begin{table}[!tb]
\begin{center}
\begin{tabular}{ |P{1.55cm}||P{1.8cm}|P{1.8cm}|P{1.60cm}|}
 \hline
 \multicolumn{4}{|c|}{\textbf{Novelty Type}} \\
 \hline
\textbf{Novelty Class} & \textbf{Beneficial} & \textbf{Detrimental} & \textbf{Irrelevant}\\
 \hline
 %level 1
Object & \begin{small}axe\end{small} & \begin{small}fence\end{small} & \begin{small}spring, arrow\end{small}\\
\hline
Attribute & \begin{small}break increase, extract-increase\end{small} & \begin{small}extract-decrease, fire-wall\end{small} & \begin{small}-\end{small}\\
\hline
Action & \begin{small}chop, jump\end{small} & \begin{small}action remap\end{small} & \begin{small}-\end{small}\\
\hline
\end{tabular}
\caption{\label{tab:novel_envs}Classification of all novelties based on the class and type.}
\end{center}
\end{table}
%

\subsection{Novel Domains}
\label{sec:novel_domains}
% In this section, we describe the novelties that we have currently implemented in the novel gridworlds.
In the present version of NovelGridworlds, we implement 12~\footnote{To prevent redundancy, we have clubbed some novelties (extract increase, decrease and arrow, spring) into one description.} exclusive novel scenarios; a detailed list is shown in Table~\ref{tab:novel_envs}. Figure~\ref{fig:venn_diag} shows the classification of each novelty based on the classes they fall under. 

\subsubsection{Fence}
\label{Fence}

This novelty falls under the category of \emph{Object} novelty. A new object called \emph{fence}, appears in the environment. The \emph{fence} encloses some or all of the existing items in the environment to obstruct the agent's access to those items. This novelty, as shown in Table~\ref{tab:novel_envs}, is classified as a detrimental novelty. In the ideal scenario, the agent should learn to \emph{break} the fence to access the items in the environment.

\subsubsection{Arrow \& Spring}

This novelty falls under the category of \emph{Object} novelty.
A new item called \emph{arrow} or \emph{spring} appears on some grid spaces in the environment. An \emph{arrow} appears in \emph{Bow-v0}, and \emph{Bow-v1}, and \emph{spring} appears in \emph{Pogostick-v0} and \emph{Pogostick-v1}. Both these novelties are classified as \emph{irrelevant} novelties. Hence, they have nothing to do with the task and the agent should learn to ignore them.

\subsubsection{Break increase}

This is an \emph{Attribute} novelty, in which the attribute of the existing item in the environment changes. In the pre-novelty scenario, if the agent breaks an object, the number of quantity of that object in the agent's inventory would be incremented by $1$.
However, in \emph{break increase} novelty, after breaking an object, the agent would receive $2$ quantities of that object in its inventory instead of $1$. This is a \emph{beneficial} novelty as the agent can collect more objects with lesser number of actions and finish the task sooner as compared to pre-novelty scenario.

\subsubsection{Extract increase/decrease}

This is an \emph{Attribute} novelty, in which the attribute of the existing item in the environment changes.
In this novelty, there is a change in the number of objects the agent gets when it executes \emph{extract} action.
This novelty can be beneficial as well as detrimental. If there is an increase in the number of objects the agent gets, it would be a beneficial novelty, whereas in case of decrease in the number of objects, it would be a \emph{detrimental} novelty.

\subsubsection{Action remapping}

This novelty falls under the class of \emph{Action} novelty, and is further classified as \emph{detrimental}. In this novelty, the effects of the actions as known to the agent are not same anymore. For example, \emph{turn-left} does not turn left anymore but make the agent move forward. We randomly shuffle the action effects with each other to remap actions among the set of available action in each task. This novelty is classified as a detrimental novelty, and the agent should learn to handle such a scenario. In a real-world scenario, this can be analogous to a system getting hacked, and the autonomous agent needs to figure out a way to get out of such a situation.

\subsubsection{Jump}

This novelty falls under the class of \emph{Action} novelty, and is further classified as \emph{beneficial}. In this novelty, a new action called \emph{jump} is added to the action space of the agent. This action makes the agent jump two blocks forward in the direction it is facing. This novelty is a beneficial novelty because it helps the agent to move faster in the environment. 

\subsubsection{Axe}

This novelty falls under two classes, \emph{Object}, and \emph{Attribute}. In this novel scenario, a new object called \emph{axe} appears in the environment, and if the agent uses it, the number of quantities received from breaking something in the environment doubles. Hence, this novelty introduces a new object and also changes the attributes of the environment. This novelty is further classified as beneficial, and ideally, the agent should learn to exploit this novelty to its advantage.

\subsubsection{Fire wall}

This novelty falls under two classes, \emph{Object}, and \emph{Attribute}. In this novel scenario, the material of the wall enclosing the arena changes to \emph{fire-wall}, and if the agent goes near it, it dies, and the episode terminates. This novel scenario is further classified as a detrimental novelty since the agent needs to be extra careful in going near the \emph{fire-wall}.

\subsubsection{Chop}

This novelty falls under two classes, \emph{Action} and \emph{Attribute}. In this novelty, a new action called \emph{chop} is added to the agent's action space. Upon using this action in front of the breakable objects in the environment, the agent gets double the number of items previously received. This novelty falls under the classification of beneficial novelty. Ideally, the agent should discover this novel action and learn to double the number of items.

\subsubsection{Axe to break}

This novelty falls under all three categories, \emph{Object}, \emph{Action}, and \emph{Attribute}. In this novelty, a new object call \emph{axe} appears in the environment, and the break action only works if the agent is holding an \emph{axe}. This novelty is detrimental to the agent's perspective since it cannot complete the task unless it discovers this novelty. This novelty is not included in the table~\ref{tab:novel_envs}, because it is falls under two categories of detrimental and beneficial and also under three classes.

\section{Experimental Results}
\label{sec:experiment}

\subsection{Agent Descriptions}
To demonstrate our environments, we ran experiments using two different types of agents on some of the base tasks and novel tasks.



\subsubsection{Planning }
\label{sec:planning_agent}
The planning agent was implemented using an extension of cognitive Distributed Integrated Affect Re-flection Cognition (DIARC) architecture~\citep{schermerhorn2006diarc, scheutz2019overview}. The planning agent parses the JSON received from the environment and generates a problem and a domain PDDL file. MetricFF~\citep{hoffmann:nebel:jair-01} planner is used to find a plan to solve the task. The plan is executed using an action executor implemented in the architecture. At each time step the agent updates its memory based on the information it receives from the environment. We used the same architecture as described in~\citep{muhammad2021novelty}. The system maintains a living, episodic and a universal memory to infer novelties in the environment. The architecture uses heuristic algorithms as proposed in~\citep{sarathy2018macgyver} to adapt to the novel situations. 

\subsubsection{Reinforcement Learning (RL) Agent}
\label{sec:rl_agent}
We created a reinforcement learning agent~\footnote{The code that trains the environment using the RL agent using stable-baselines is available at ~\url{https://github.com/goelshivam1210/adaptive_agent}. It also has the necessary PDDL files and other data to reproduce results presented in the paper.} with LIDAR sensors. The sensors send $8$ beams at $45$ degree increments to each reachable object in the environment. Each beam returns a euclidean distance from the agent's current location to the object. The agent also has a vector representing the current inventory. 
For the sake of simplicity, the \emph{crafting-table} was made an unbreakable object.
The reward function was also shaped to make learning easier. A reward of $+10$ was given if the agent collected one of the ingredients for the task completion. A reward of $+50$ was given for achieving the goal, and $-1$ for each time step otherwise.
For training the RL agent, we used the Proximal Policy Optimization (PPO2)~\citep{schulman2017proximal} implementation in the open-source
Stable Baselines package~\citep{stable-baselines3}, with default hyperparameters. We chose an off the shelf implemented RL algorithm to demonstrate how users can easily plug in existing algorithms with our environment.
% A regular policy gradient~\citep{sutton1999policy} with $\epsilon$ greedy policy selection was used to train with decaying exploration constant $\epsilon$ beginning at 0.1 and decaying towards $0.05$.  $\epsilon$ is decayed exponentially using the formula $\epsilon(t)=$ $\epsilon_{\min }+\left(\epsilon_{\max }-\epsilon_{\min }\right) e^{-\lambda t},$ where $\lambda=-\log (0.01) / N, N$ being the maximum number of episodes.

\subsection{Evaluation Metrics}
\label{sec:eval_metrics}



% In this section, we describe the various novelties which are implemented in the repository we publish. We also describe the various agents we have developed to solve some of these novelties for the users to get an idea on how to plug and play. We present the agents in a separate repository which is also available online.
% We run experiments on $4$ environments as mentioned in Section~\ref{sec:environment}. Each of these 4 environment have some novelties implemented as listed in the table\ref{tab:novel_envs}.
% It is pertinent for the novelties to overlap between classes and its hard to put them in classes strictly belonging to one. Hence, we also categorize each novel environment as shown in~\ref{fig:venn_diag}.
In this section, we describe our proposed evaluation metrics for measuring performance of novelty-centric agents in both planning systems and learning systems.

\subsubsection{Base task performance metrics} 
These metrics evaluate the performance of the agents on the base task (pre-novelty scenario). In both the planning and reinforcement learning agents, we measure the number of time steps the agent takes to complete the base task. Specifically for the RL agent, we also measure the percentage of times the agent achieves the final goal. In the case of the planning agent, the agent always succeeds in reaching the goal in the pre-novelty environment.

\subsubsection{Novelty metrics} 
\label{sec:metrics}
These metrics evaluate the performance of the agent after the novelty is introduced.
The metrics are reliant on two criterion, namely detection and adaption. Hence we define three metrics to evaluate the performance of the agent in the novelty scenario.
%
\begin{itemize}
    % \item Metric M1 (Detection Performance): It is the average of precision and recall of the number of times the agent detects a novelty in the environment.
    \item Metric M1 (Reaction Performance): This metric is defined as the ratio of the cumulative action cost that the agent accumulates in the post-novelty scenario to the action cost accumulated in the pre-novelty scenario. In other words, if the ratio is less than $1$, the performance is improved; if the ratio is greater than $1$, the performance is reduced and if the ratio is $1$, the performance is same.
    \item Metric M2 (Steps): This metric records the number of steps taken by the agent to solve the task in the post-novelty scenario.
    \item Metric M3 (Detection Rate): This metric records the binary value of whether the agent reports the novelty or not. This metric is only specific to the planning architecture.
\end{itemize}
%
To demonstrate the RL agent's performance, we plot the learning curves. We divide the learning curve into two sections; section 1 presents the agent's learning performance in the pre-novelty scenario and section 2 represents the agent's performance in the post-novelty scenario. We present the test results of the agent over a course of $10$ trials. In each trial, we save $20$ models uniformly distributed across the training regime. We perform a test on each model, by running it on the corresponding environment for $25$ episodes, each with a maximum of $200$ steps. We compute the cumulative reward in every episode and report the cumulative mean reward of all the $25$ episodes. This score averaged across $10$ trials is plotted against the number of time steps the agent is trained. 
To do fair comparison to actually know whether the RL agent learned to solve the task, we also report the success rate of achieving the goal task. To compute the success rate, we run $25$ tests on the final models of each training regime and report the percentage of times the agent successfully reaches the goal state.

\subsection{Results}
\label{sec:results}
\subsubsection{Planning Agent}
We evaluated the planning agent on all the $4$ tasks, and $3$ novelty scenarios in each task.
\begin{enumerate}
    \item Fence
    \item Arrow/Spring
    \item Axe
    % \item Axe to break
\end{enumerate}

Table~\ref{tab:results_planning} shows the performance of the planning agent. 
We ran the agent for $10$ episodes. In each episode the agent was given the environment dynamics, and a goal to solve. Keeping the goal same across the episodes, we introduce the desired novelty in the $6^{\text{th}}$ episode. After the $6^{\text{th}}$ episode, all the subsequent episodes were loaded with the same novel environment. At each time step in an episode, the agent gets a step cost for the action taken from the environment. The cumulative step cost is used to compute the metric M1. The agent was able to achieve the goal in all the episodes. The table shows the mean and standard deviation of five non-novel episodes and five novel episodes.
The planning agent we evaluate achieved a $100 \%$ detection rate in all the novelty scenarios. Hence, the metric M3 for the planning agent was a perfect score.
%
\begin{table}[!tb]
\begin{center}
\begin{tabular}{|P{1.0cm}||P{1.4cm}|P{1.3cm}|P{1.6cm}|P{1.2cm}|}
 \hline
 \multicolumn{5}{|c|}{\textbf{NovelGridworld-Bow-v0}} \\
 \hline
\textbf{Metrics} & \begin{small}\textbf{Default}\end{small} & \begin{small}\textbf{Fence}\end{small} & \begin{small}\textbf{Arrow/Spring}\end{small} &\begin{small}\textbf{Axe}\end{small}\\
 \hline
 %level 1
M1 & \begin{small}$1.00 \pm 0.00$ \end{small} & \begin{small}$1.80 \pm 0.15$\end{small} & \begin{small}$1.11\pm0.135$\end{small} & \begin{small} $1.02\pm0.04$\end{small}\\
M2 & \begin{small}$87.6\pm16.8$\end{small} & \begin{small}$123.0\pm18.0$\end{small} & \begin{small}$90.6\pm15.1$\end{small} & \begin{small}$98.0\pm15.9$ \end{small}\\
% M3 & \begin{small}\end{small} & \begin{small}\end{small} & \begin{small}\end{small} & \begin{small} \end{small}& \begin{small} \end{small}\\
% M4 & \begin{small}\end{small} & \begin{small}\end{small} & \begin{small}\end{small} & \begin{small} \end{small}\\
 \hline
 \multicolumn{5}{|c|}{\textbf{NovelGridworld-Bow-v1}} \\
 \hline
 %level 1
M1 & \begin{small}$1.00\pm0.00$\end{small} & \begin{small}$1.82\pm0.24$\end{small} & \begin{small}$1.00\pm0.02$\end{small} & \begin{small}$0.99\pm0.02$\end{small}\\
M2 & \begin{small}$36.8\pm6.5$\end{small} & \begin{small}$54.2\pm4.8$\end{small} & \begin{small}$37\pm5.4$\end{small} & \begin{small}$35.4\pm4.0$\end{small}\\
% M3 & \begin{small}\end{small} & \begin{small}\end{small} & \begin{small}\end{small} & \begin{small} \end{small}& \begin{small} \end{small}\\
% M4 & \begin{small}\end{small} & \begin{small}\end{small} & \begin{small}\end{small} & \begin{small} \end{small}\\
 \hline
 \multicolumn{5}{|c|}{\textbf{NovelGridworld-Pogostick-v0}} \\
 \hline
 %level 1
M1 & \begin{small}$1.00\pm0.00$\end{small} & \begin{small}$1.16\pm0.08$\end{small} & \begin{small}$0.97\pm0.03$\end{small} & \begin{small}$0.99\pm0.04$ \end{small}\\
M2 & \begin{small}$127.8\pm9.0$\end{small} & \begin{small}$122.0\pm19.8$\end{small} & \begin{small}$105.6\pm15.7$\end{small} & \begin{small}$114.8\pm4.0$\end{small}\\
% M3 & \begin{small}\end{small} & \begin{small}\end{small} & \begin{small}\end{small} & \begin{small} \end{small}& \begin{small} \end{small}\\
% M4 & \begin{small}\end{small} & \begin{small}\end{small} & \begin{small}\end{small} & \begin{small} \end{small}\\
 \hline
 \multicolumn{5}{|c|}{\textbf{NovelGridworld-Pogostick-v1}} \\
  \hline
 %level 1
M1 & \begin{small}$1.00\pm0.00$\end{small} & \begin{small}$1.34\pm0.06$\end{small} & \begin{small}$1.07\pm0.05$\end{small} & \begin{small}$0.99\pm0.02$\end{small}\\
M2 & \begin{small}$94.2\pm12.9$\end{small} & \begin{small}$105.2\pm14.5$\end{small} & \begin{small}$100.8\pm22.58$\end{small} & \begin{small}$85.0\pm15.4$\end{small}\\
% M3 & \begin{small}\end{small} & \begin{small}\end{small} & \begin{small}\end{small} & \begin{small} \end{small}& \begin{small} \end{small}\\
% M4 & \begin{small}\end{small} & \begin{small}\end{small} & \begin{small}\end{small} & \begin{small} \end{small}\\
 \hline
\end{tabular}
\caption{\label{tab:results_planning}~Table shows the results of the performance of the planning agent on all the four tasks. In each task, it shows pre-novelty (Default) scenario and 3 post-novelty (Fence, Arrow/Spring, Axe) scenarios.}
\end{center}
\end{table}
%

It is important to note that the metric M1 for fence novelty in all the environments is greater than 1, which clearly shows that the agent needs to spend a higher cost to break the fences to achieve the goal. Since the fence is a detrimental novelty, the agent accumulates more cost while handling the novelty. The metric M2, in this case also shows that the agent needs more steps than the default case to finish the task.

It can be seen from Table~\ref{tab:results_planning}, that the agent seeks to improve its performance in the \emph{axe} novelty. However, the improvement is negligible. A possible reason for negligible improvement can be that the agent takes more steps to complete the tasks in novel cases (since it needs to navigate to the axe and hold it). Hence, it incurs more step cost than it improves by using the axe. Moreover, we noticed that in some of the episodes the agent did not even use the axe to benefit from using it. In the case of the irrelevant novelty (arrow/spring), the agent's performance is nearly the same as the default case, as expected. 

% We notice that the metric M2 for the pogostick-v0 shows a decrease in number of steps when adding a novelty. A reason for that behavior can be that the plan generated is sub-optimal, and it does not account action costs in planning, due to which it takes less steps in a complex scenario.

\subsubsection{Reinforcement learning agent}

The RL agent was evaluated on the \emph{NovelGridworld-Bow-v0} task. We evaluate the performance on $3$ novelty scenarios in this task. The agent was evaluated on these novelties for demonstration purposes and users are encouraged to develop methods to handle more complex novelties as implemented in the environment.
\begin{enumerate}
    \item Action remapping 
    \item Fire wall
    \item Break increase
\end{enumerate}

%
\begin{figure}[!tb]
  \centering
  \includegraphics[width=1\linewidth]{results_latest.png}
  \caption{Figure shows the learning performance of the RL agent on the task of \emph{NovelGridworld-Bow-v0}. The pink line in the center (2 Million) separates the performance of the agent in the pre-novelty scenario to the agent's performance in cases of novelty.}
  \label{fig:rl_results_plot}
\end{figure}
%
Figure~\ref{fig:rl_results_plot} shows the performance of the RL agent on the base task (shown in red) and in the three novelty scenarios, namely \emph{breakincrease} (shown in orange), \emph{fire-wall} (shown in green), and \emph{action remapping} (shown in blue). The figure was plotted as described earlier in Section~\ref{sec:eval_metrics} .
It can be seen from Figure~\ref{fig:rl_results_plot}, that the RL agent performs well in solving the task after being trained for around 1.8 million time steps. Row 1 of Table~\ref{tab:results_rl}, shows the success rate (percentage of times the agent reaches the goal state) of the default task after 1.8 million time steps of training.
Its encouraging to see that the agent improves its performance in the post novelty scenario in the case of \emph{breakincrease} novelty (orange line in Figure ~\ref{fig:rl_results_plot}). This was expected since it is a beneficial novelty. However, if we notice in Table~\ref{tab:results_rl}, the success rate of the agent drops considerably. One of the possible reasons for increase in the total rewards but decrease in the performance can be that the agent tries to collect resources and hoard reward without completing the final task.
We also notice a sharp decline in agent's performance when the \emph{fire-wall} novelty is introduced. \emph{Fire-wall} novelty is a detrimental novelty and the task becomes quite hard to solve, hence we see that the agent never recovers in this novel scenario. Table~\ref{tab:results_rl}, also reflects the poor performance of the agent in the \emph{fire-wall} novelty.

Action remapping novelty however, showed almost similar performance when compared to the pre-novelty performance. Provided that its a \emph{detrimental} novelty, the behavior is surprising. 
A possible reason for this behavior can be that the action space is too small (5 discrete actions), and the actions are randomly shuffled. Therefore, in some trials the agent might have got lucky and only $2$ actions might have inter-changed. We discovered this, when we inspected the agent's performance on each trial, and found out that it dipped considerably (to almost -150 in the reward curve), when all the actions were remapped. However, it did bounce back to original performances in around $200,000$ environment interactions. Table~\ref{tab:results_rl}, however, shows the decline of $12\%$ in the performance of the agent.

The RL agent was not tailored explicitly to handle novelties. Hence, we could not evaluate it on all novelties. However, RL agents do not have access to the transition dynamics and learn from interactions, hence some level of non-staionarity such as, action remapping and break increase was handled by the RL agent.
The addition of a new item in the agent's state space is particularly hard, especially when the state space representation of the RL agent is symbolic. Hence, we did not evaluate the agent on \emph{object} novelties, and leave it as an open challenge for the RL community

%
\begin{table}
\begin{center}
\begin{tabular}{ |P{3cm}||P{2cm}| }
 \hline
 \multicolumn{2}{|c|}{\textbf{NovelGridworld-Bow-v0}} \\
 \hline
 \textbf{Novelty}& \textbf{Success (\%)} \\
 \hline
 Default   & $92.0 \pm 20.8$ \\
  Action remapping   & $80.00 \pm 40.00$ \\
   Fire wall   & $0.00 \pm 0.00$ \\
    Break increase & $29.2 \pm 19.6$ \\
 \hline
\end{tabular}
\caption{\label{tab:results_rl}Table shows the results of the performance of the reinforcement learning agent on the base task of \emph{NovelGridworld-Bow-v0} and three novelty scenarios. The values denote the mean and standard deviations of the success rate of solving the final task. The values were calculated at 1.8 M timestep for the default case and 4 M time step for all the novelty scenarios.}
\end{center}
\end{table}
%

\section{Discussion}
\label{sec:discussion}

The architecture of NovelGridworlds is designed for cross-platform applicability and cross-agent compatibility. As shown in Section ~\ref{sec:experiment}, NovelGridworlds is capable of prototyping algorithms both in the planning domains and machine learning domains. It is a fast and lightweight simulation with a socket connection~\footnote{We have detailed out all the special features and how to use them in the readme of our environment which can be found at this link: \url{https://github.com/gtatiya/gym-novel-gridworlds/blob/master/README.md}} for communication between the agent and the environment. Support for socket connection provides an efficient way to connect the environment with various agents, implemented using different programming languages.
We demonstrate this ability by evaluating the environment using agents developed in Java as well as Python.

Another striking feature of the NovelGridworlds is the verbose information about the state of the world. This information can be exploited by Natural Language Processing (NLP) algorithms to perform novelty reasoning. Hence, combining NLP algorithms with the RL and Planning frameworks could be an important research direction that can be pursued using our environment. Our framework is also flexible in adding multiple novelties to the same task, thereby making the task harder and generating a curriculum. This functionality provides the capability of training a curriculum for reinforcement learning agents~\citep{narvekar2020curriculum}.

% One of the primary goals of building this infrastructure is to provide a platform for developing algorithms dealing with novelties in open worlds. We provide a rich taxonomy of the novelties in this paper (Section~\ref{sec:novel_domains}), and we have implemented novelty wrappers in the code-base for easy access to the novel domains.
Reinforcement learning algorithms deal with sample inefficiency, and many researchers work in the direction of pre-training using human demonstrations~\citep{de2019pre}. To that extent, we provide keyboard interface support in our environment and provide support for recording trajectories, which can be further used to perform learning from demonstrations. We also provide code examples of wrappers for limiting action space, implementing various observation spaces essential in RL algorithms training. 
% \subsubsection{JSON interface}
% \subsubsection{Keyboard interface}
% \subsubsection{Socket interface}
% \subsubsection{Novelty taxonomy}
% \subsubsection{}
% In this section talk about the capabilities of novel gridworld.
% As shown in section 

% - Discuss the usefulness of this environment.
% - How people can use it to develop algorithms for novel scenarios.\\
% - How a combination of planning and RL is easy to develop since the environment is friendly to both the learning methods.\\
% - How easy it is to generate novel environments. \\
% - keyboard interface \\
% - Socket connection capability for cross platform adaptability.\\
% - Curriculum generation tasks are quite suitable for this type of domain. \\
% - Crafting tasks.. \\



\section{Conclusion \& Future Work}
\label{sec:conclusion}
We presented {\em NovelGridworlds}, a robust open-source benchmarking environment for developing and evaluating agents capable of handling novelties in open worlds, together with a rich taxonomy of how novelties in this domain are characterized.  Results from running and evaluating both planning and learning agents in different tasks show the utility of {\em NovelGridworlds} for comparing and quantifying different approaches to open-world AI.

Moving forward, we plan to include multi-agent support that will allow for training and benchmarking multi-agent RL and planning algorithms.  And while the current version only allows for static novelties (i.e., novelties that remain the same throughout a task run), future versions of the system will also include dynamic novelties that can change at any time during task performance.  This will enable even more challenging evaluation settings that can test an agent's ability to quickly react to novelties.

%An example of a dynamic novelty can be in terms of addition of fences in the environment to block the %crafting table, after the agent completes a specific part of the task. The agent in this case needs to %quickly adapt to that change in its representation of the world. 


% - More tasks and more novelties \\

% We plan to improve many things as a part of future work. An integral part of dealing with open-world novelties is adapting algorithms to the robotics domain. Therefore, we plan to add more tasks that are relevant to robots and continuous control. We also plan to include multi-agent support to train and benchmark Multi-agent RL algorithms. The current version of the NovelGridworlds does not support visual information; as a crucial future work, we plan to integrate rich visual output. Visual information can help researchers perform novelty detection and adaption using computer vision and deep learning techniques. In the current version, we deal with static novelties; in other words, novelties are the same throughout the episode. However, as a part of future work, we plan to include dynamic novelties, in which the novelties are introduced during the episode.
% - multi-agent support \\
% As a part of future work we also plan to integrate...???
% We also ..... \\
% - Automatically generating PDDLs is an important task in dealing with novelties in open world specially in planning domains. 
% - 










%  \While{True}{
%   instructions\;
%   re-parameterize the sub-policy $\mu_{t}^'$ \;
%   //$\nu_t$ is the new option/sub-policy  \;
%   $\nu_{t} = \mu_{t}'$\;
%   \While {reach goal sub-policy $\nu_{t}^{*}$}{
%     until steps complete\;
%     reach goal  \;
%   }
%  }
% You can typeset all sorts of in-line mathematical expressions 
% with the usual \verb|$...$| construct, as in 
% $\Diamond\Diamond\varphi \rightarrow \Diamond\varphi$ or 
% $\boldsymbol{R} = (R_1,\ldots,R_n)$.
% For more complex expressions, it may often be preferable to use one of
% the various equation-type environments available in \LaTeX\@, as shown 
% in the following example:

% Here is a second example for an equation: 
% %
% \begin{eqnarray}\label{eq:vcg}
% p_i(\boldsymbol{\hat{v}}) & = &
% \sum_{j \neq i} \hat{v}_j(f(\boldsymbol{\hat{v}}_{-i})) - 
% \sum_{j \neq i} \hat{v}_j(f(\boldsymbol{\hat{v}})) 
% \end{eqnarray}
% %
% Use the usual combination of `\verb|\label|' and `\verb|\ref|' to refer
% to numbered equations, such as Equation~(\ref{eq:vcg}) above. Of course,
% introducing numbers in the first place is only helpful if you in fact 
% need to refer back to the equation in question elsewhere in the paper.


% \subsection{Algorithm}

% Use the `\texttt{table}' environment (or its variant `\texttt{table*}')
% in combination with the `\texttt{tabular}' environment to typeset tables
% as floating objects. The `\texttt{aamas}' document class includes the 
% `\texttt{booktabs}' package for preparing high-quality tables. Tables 
% are often placed at the top of a page near their initial cite, as done 
% here for Table~\ref{tab:locations}.

% % \begin{table}[t]
% %   \caption{Locations of the first five editions of AAMAS}
% %   \label{tab:locations}
% %   \begin{tabular}{rll}\toprule
% %     \textit{Year} & \textit{City} & \textit{Country} \\ \midrule
% %     2002 & Bologna & Italy \\
% %     2003 & Melbourne & Australia \\
% %     2004 & New York City & USA \\
% %     2005 & Utrecht & The Netherlands \\
% %     2006 & Hakodate & Japan \\ \bottomrule
% %   \end{tabular}
% % \end{table}

% The caption of a table should be placed \emph{above} the table. 
% Always use the `\verb|\midrule|' command to separate header rows from 
% data rows, and use it only for this purpose. This enables assistive 
% technologies to recognise table headers and support their users in 
% navigating tables more easily.

% %%% The following command should be issued somewhere in the first column 
% %%% of the final page of your paper.
% \balance

% Use the `\texttt{figure}' environment for figures. If your figure 
% contains third-party material, make sure to clearly identify it as such.
% Every figure should include a caption, and this caption should be placed 
% \emph{below} the figure itself, as shown here for Figure~\ref{fig:logo}.

% % \begin{figure}[h]
% %   \centering
% %   \includegraphics[width=\linewidth]{aamas21logo}
% %   \caption{The logo of AAMAS-2021}
% %   \label{fig:logo}
% %   \Description{Logo showing the words "AAMAS 2021, London, UK", 
% %   with the 0 resembling the London Eye and the 1 resembling Big Ben.}
% % \end{figure}

% In addition, every figure should also have a figure description, unless
% it is purely decorative. Use the `\verb|\Description|' command for this 
% purpose. These descriptions will not be printed but can be used to 
% convey what's in an image to someone who cannot see it. They are also 
% used by search engine crawlers for indexing images, and when images 
% cannot be loaded. A figure description must consist of unformatted plain 
% text of up to 2000~characters. For example, the definition of 
% Figure~\ref{fig:logo} in the source file of this document includes the 
% following description: ``Logo showing the words `AAMAS 2021, London, 
% UK', with the 0 resembling the London Eye and the 1 resembling Big 
% Ben.'' For more information on how best to write figure descriptions 
% and why doing so is important, consult the information available here: 
% %
% \begin{center}
% \url{https://www.acm.org/publications/taps/describing-figures/}
% \end{center}
% %
% The use of colour in figures and graphs is permitted, provided they 
% remain readable when printed in greyscale and provided they are 
% intelligible also for people with a colour vision deficiency.

%%%%%%%%%%%%%%%%%%%%%%%%%%%%%%%%%%%%%%%%%%%%%%%%%%%%%%%%%%%%%%%%%%%%%%%%

% \section{Citations and References}
  
% The use of the \BibTeX\ to prepare your list of references is highly 
% recommended. To include the references at the end of your document, put 
% the following two commands just before the `\verb|\end{document}|' 
% command in your source file:
% %
% \begin{verbatim}
%   \bibliographystyle{ACM-Reference-Format}
%   \bibliography{mybibfile}
% \end{verbatim}
% %
% Here we assume that `\texttt{mybibfile.bib}' is the name of your 
% \BibTeX\ file. Use the `\verb|\cite|' command to produce citations 
% to your references. Here are a few examples for citations of journal 
% articles~\cite{GrKr96,WoJe95}, books~\cite{Knu97}, articles in 
% conference proceedings~\cite{Hag1993}, technical reports~\cite{Har78},
% Master's and PhD theses~\cite{Ani03,Cla85}, online videos~\cite{Oba08}, 
% datasets~\cite{AnMC13}, and patents~\cite{Sci09}. Both citations and 
% references are numbered by default. 

% Make sure you provide complete and correct bibliographic information 
% for all your references, and list authors with their full names 
% (``Donald E.\ Knuth'') rather than just initials (``D.\ E.\ Knuth''). 

%%%%%%%%%%%%%%%%%%%%%%%%%%%%%%%%%%%%%%%%%%%%%%%%%%%%%%%%%%%%%%%%%%%%%%%%

%%% The acknowledgments section is defined using the "acks" environment
%%% (rather than an unnumbered section). The use of this environment 
%%% ensures the proper identification of the section in the article 
%%% metadata as well as the consistent spelling of the heading.


% MS: acks always come after acceptance to not bias reviewers
%
\begin{acks}
This work was funded in part by DARPA grant W911NF-20-2-0006.
%Add the DARPA funding this project is associated with, and the people involved in the SAIL-ON team at %Tufts.
%If you wish to include any acknowledgments in your paper (e.g., to 
%people or funding agencies), please do so using the `\texttt{acks}' 
%environment. Note that the text of your acknowledgments will be omitted
%if you compile your document with the `\texttt{anonymous}' option.
\end{acks}

%%%%%%%%%%%%%%%%%%%%%%%%%%%%%%%%%%%%%%%%%%%%%%%%%%%%%%%%%%%%%%%%%%%%%%%%

%%% The next two lines define, first, the bibliography style to be 
%%% applied, and, second, the bibliography file to be used.

\bibliographystyle{ACM-Reference-Format} 
\bibliography{sample}

%%%%%%%%%%%%%%%%%%%%%%%%%%%%%%%%%%%%%%%%%%%%%%%%%%%%%%%%%%%%%%%%%%%%%%%%

\end{document}

%%%%%%%%%%%%%%%%%%%%%%%%%%%%%%%%%%%%%%%%%%%%%%%%%%%%%%%%%%%%%%%%%%%%%%%%


